% applications/examples-DE.tex
\chapter{Anwendungsbeispiele}
\label{chap:applications-examples}

Dieses Kapitel veranschaulicht repräsentative Anwendungsmuster, die
etablierte AIM-Begriffe verwenden.

\section{Beispielmuster A: Vergleich von Übergängen}

Zwei Übergangsgesetze können identische Randbedingungen erfüllen, dabei
aber unterschiedliche innere Krümmungsverläufe und folglich
unterschiedliche laufzeitabgeleitete dynamische Größen erzeugen.

Ein kanonischer Indikator bleibt
\[
a_{\mathrm{lat}}(s)=v^2\kappa(s),
\]
womit die betriebliche Konsequenz sowohl von der Geometrie als auch vom
Geschwindigkeitskontext abhängt.

\section{Beispielmuster B: Überhöhungsgekoppelte Interpretation}

Die Anwendungsinterpretation verwendet gekoppelte Krümmungs- und
Überhöhungsgrößen sowie Laufzeitoperatoren, um Folgen für Zulässigkeit und
Komfort zu bewerten. Abgeleitete Größen bleiben nachgelagerte
Laufzeitergebnisse.

\section{Beispielmuster C: Vergleich realisierter Zustände}

Soll- und realisierte Zustände werden in realisierungsbewussten Abläufen
verglichen. Das Beispiel richtet sich auf die Ingenieurkonsequenz
(Instandhaltung oder betriebliche Anpassung), nicht auf eine Neudefinition
vorgelagerter Semantik.

\section{Verbindung zu AIM}
\begin{summary}
Die Beispiele zeigen, wie kanonische AIM-Begriffe in Ingenieurkontexten
für Vergleich, Auswertung und Entscheidungsunterstützung verwendet werden.

Sie beleuchten bewusst die etablierte Architektur, anstatt neue
Begriffsebenen einzuführen.
\end{summary}
